% GENERATED FILE. Edit canonical lessons, then run npm run generate:books.
% canonical-source-hash: fnv1a64:33a91364c507c486
% canonical-lessons: RU-C251-byk, RU-C251-petukh, RU-C251-gus, RU-C251-golub, RU-C251-vorobey

\chapter{Bull, Cockerel, Goose, Pigeon, Sparrow}
\label{ch:ru-bull-cockerel-goose-pigeon-sparrow}
\hlchaptermodality{251}

\begin{chapteropening}
\textbf{By the end of this chapter:} \emph{I can say a bull, a cockerel, a goose, a pigeon and a sparrow.}

Complete the last lesson of chapter 251: I can say a bull, a cockerel, a goose, a pigeon and a sparrow.
\end{chapteropening}

\section[\texorpdfstring{\ru{бык} (byk)}{бык (byk)}]{\ru{бык} (byk) — a bull}

\label{lesson:RU-C251-byk}



\begin{quote}
Before the new one: say the Russian for a cake, then the Russian for a sweet.
\end{quote}

\subsection*{You'll want to know: \ru{бык}}
\textbf{\ru{бык}} (\emph{byk}) — \textquotedblleft{}a bull\textquotedblright{}.

\begin{grammarlens}[title={What you've built}]
One more word for this chapter.
\end{grammarlens}

\subsection*{Your turn}
\begin{itemize}
\raggedright
  \item \emph{Say it:} \emph{byk}
  \item \emph{Say it:} \emph{byk}, once more
  \item \emph{Say it:} say \emph{byk}
  \item \emph{Recall:} say the Russian for a cake, then the Russian for a sweet, then say \emph{byk} again
\end{itemize}

\begin{tcolorbox}[breakable,colback=teal!4,colframe=teal!35!black,title={Before you move on}]
What does \ru{бык} mean? (\textquotedblleft{}A bull\textquotedblright{}.) Say it once more.
\end{tcolorbox}
\section[\texorpdfstring{\ru{петух} (petúkh)}{петух (petúkh)}]{\ru{петух} (petúkh) — a cockerel}

\label{lesson:RU-C251-petukh}



\begin{quote}
Before the new one: say the Russian for a sweet, then the Russian for a bull.
\end{quote}

\subsection*{You'll want to know: \ru{петух}}
\textbf{\ru{петух}} (\emph{petúkh}) — \textquotedblleft{}a cockerel\textquotedblright{}.

\begin{grammarlens}[title={What you've built}]
One more word for this chapter.
\end{grammarlens}

\subsection*{Your turn}
\begin{itemize}
\raggedright
  \item \emph{Say it:} \emph{petúkh}
  \item \emph{Say it:} \emph{petúkh}, once more
  \item \emph{Say it:} say \emph{petúkh}
  \item \emph{Recall:} say the Russian for a sweet, then the Russian for a bull, then say \emph{petúkh} again
\end{itemize}

\begin{tcolorbox}[breakable,colback=teal!4,colframe=teal!35!black,title={Before you move on}]
What does \ru{петух} mean? (\textquotedblleft{}A cockerel\textquotedblright{}.) Say it once more.
\end{tcolorbox}
\section[\texorpdfstring{\ru{гусь} (gus)}{гусь (gus)}]{\ru{гусь} (gus) — a goose}

\label{lesson:RU-C251-gus}



\begin{quote}
Before the new one: say the Russian for a bull, then the Russian for a cockerel.
\end{quote}

\subsection*{You'll want to know: \ru{гусь}}
\textbf{\ru{гусь}} (\emph{gus}) — \textquotedblleft{}a goose\textquotedblright{}.

\begin{grammarlens}[title={What you've built}]
One more word for this chapter.
\end{grammarlens}

\subsection*{Your turn}
\begin{itemize}
\raggedright
  \item \emph{Say it:} \emph{gus}
  \item \emph{Say it:} \emph{gus}, once more
  \item \emph{Say it:} say \emph{gus}
  \item \emph{Recall:} say the Russian for a bull, then the Russian for a cockerel, then say \emph{gus} again
\end{itemize}

\begin{tcolorbox}[breakable,colback=teal!4,colframe=teal!35!black,title={Before you move on}]
What does \ru{гусь} mean? (\textquotedblleft{}A goose\textquotedblright{}.) Say it once more.
\end{tcolorbox}
\section[\texorpdfstring{\ru{голубь} (gólub)}{голубь (gólub)}]{\ru{голубь} (gólub) — a pigeon}

\label{lesson:RU-C251-golub}



\begin{quote}
Before the new one: say the Russian for a cockerel, then the Russian for a goose.
\end{quote}

\subsection*{You'll want to know: \ru{голубь}}
\textbf{\ru{голубь}} (\emph{gólub}) — \textquotedblleft{}a pigeon\textquotedblright{}.

\begin{grammarlens}[title={What you've built}]
One more word for this chapter.
\end{grammarlens}

\subsection*{Your turn}
\begin{itemize}
\raggedright
  \item \emph{Say it:} \emph{gólub}
  \item \emph{Say it:} \emph{gólub}, once more
  \item \emph{Say it:} say \emph{gólub}
  \item \emph{Recall:} say the Russian for a cockerel, then the Russian for a goose, then say \emph{gólub} again
\end{itemize}

\begin{tcolorbox}[breakable,colback=teal!4,colframe=teal!35!black,title={Before you move on}]
What does \ru{голубь} mean? (\textquotedblleft{}A pigeon\textquotedblright{}.) Say it once more.
\end{tcolorbox}
\section[\texorpdfstring{\ru{воробей} (vorobéy)}{воробей (vorobéy)}]{\ru{воробей} (vorobéy) — a sparrow}

\label{lesson:RU-C251-vorobey}



\begin{quote}
Before the new one: say the Russian for a goose, then the Russian for a pigeon.
\end{quote}

\subsection*{You'll want to know: \ru{воробей}}
\textbf{\ru{воробей}} (\emph{vorobéy}) — \textquotedblleft{}a sparrow\textquotedblright{}.

\begin{grammarlens}[title={What you've built}]
One more word for this chapter.
\end{grammarlens}

\subsection*{Your turn}
\begin{itemize}
\raggedright
  \item \emph{Say it:} \emph{vorobéy}
  \item \emph{Say it:} \emph{vorobéy}, once more
  \item \emph{Say it:} say \emph{vorobéy}
  \item \emph{Recall:} say the Russian for a goose, then the Russian for a pigeon, then say \emph{vorobéy} again
\end{itemize}

\begin{tcolorbox}[breakable,colback=teal!4,colframe=teal!35!black,title={Before you move on}]
What does \ru{воробей} mean? (\textquotedblleft{}A sparrow\textquotedblright{}.) Say it once more.
\end{tcolorbox}
